\section{Discussion}

% =====================================================================
% [YOURS] — the discussion paragraphs are argument and a reviewer will
% probe them. Spine and facts below; write the prose yourself.
% =====================================================================
%
% PARAGRAPH 1 — what the sign-flip means for the field.
%   - Published FM-vs-supervised attributions that rest on one comparator
%     architecture are not safe to read as measurements.
%   - The fix is cheap: run the control on more than one architecture and
%     report the raw gaps.
%   - Generalises past EEG to any pretrained-vs-task-specific comparison
%     where the two arms are fed different preprocessing.
%
% PARAGRAPH 2 — what we are NOT saying.
%   - Not "foundation models do not work for EEG." Two models, two datasets,
%     one task family, and a clear 2-class counterexample in Section V-D.
%   - Not a claim about pretraining in general; see Section IV-C on why the
%     random-init control is not ours.
%
% PARAGRAPH 3 — one sentence only on future work. Name edge deployment and
% quantization-robust calibration as the open question. DO NOT promise
% hardware in a paper that does not have it.
% (docs/unified_project_and_freeze.md S7 rule 5.)


% =====================================================================
% RAW MATERIAL, 2026-09-20. REWRITE IN YOUR OWN VOICE BEFORE SUBMITTING.
% manuscript_s3_s4_spec.md names the discussion as "what a judge probes."
% =====================================================================

The result in Section~\ref{sec:decomp} has a consequence beyond the two
encoders audited here. Where a published comparison between a pretrained model
and a task-specific one attributes part of the gap to preprocessing using a
single comparator architecture, our data suggest that attribution should be
read with caution: the same pipeline helped one architecture by 0.078 accuracy
and cost another 0.088, so the choice of comparator alone spans a large share
of the quantity being apportioned, and none of the individual terms was
separately established at $n=9$. The remedy is cheap. Run the
control on more than one architecture, and report the raw gaps rather than a
percentage split, because a share is only interpretable when its denominator
has a stable sign. Nothing about this argument is specific to EEG; it applies
wherever a pretrained model and a task-specific baseline are necessarily fed
different inputs.

Two things this paper does not claim. It is not evidence that pretrained
% [UNSCOPE] "two datasets" -> three once results_n54.tex lands.
encoders fail on EEG generally: we audit two models on two datasets within one
task family, and report a two-class setting in which a fine-tuned CBraMod is
competitive with every supervised arm. Nor is the architecture-matched
random-initialisation control ours --- it is established practice in recent
audits~\cite{kontras2026neuroatlas,banville2026neuralbench,zare2026negativecontrol}.
What we add is the measurement of the input-matching term and the observation
that it is not stable.

The open question we consider most useful is deployment. Calibration here is
restored by a scalar fitted on validation data in floating point, and whether
that scalar and the abstention thresholds derived from it survive the reduced
numeric precision of on-device inference is untested, in this work and, to our
knowledge, elsewhere.
